\section{Limitations}
\label{sec:limitations}

\begin{table}[htbp]
\centering
\caption{Summary of model and evidentiary limitations.}
\label{tab:limitations}
\adjustbox{max width=\textwidth}{\begin{tabular}{@{}p{0.32\textwidth}p{0.63\textwidth}@{}}
\toprule
\textbf{Limitation} & \textbf{Scope / disclosed status} \\
\midrule
Parameter domain \& duration & Results conditional on tested mechanism set, parameter ranges, and run length ($\le$700 steps); not a claim about untested regions. \\
Collapse still modal & 56.6\%--78.1\% collapse across L1.2b pooled samples even after stratification; non-collapse variation is the minority of the sampled space. \\
\texttt{shared\_lineage\_persistence} unevaluated & Would require a 4th 576-run re-simulation with per-run lineage export; judged low-value given the continuum/null-replication result; disclosed, not filled here. \\
Temporal-representation validation subset & 4 burn-in definitions validated on 12 of 576 discovery runs, for compute-budget reasons stated in advance. \\
Supplementary significance tests & L1.2 Section 13 Wilcoxon tests ($n$=8 pairs) are supplementary; L1.2b's 13-factor decomposition reports effect sizes, not significance tests, by design. \\
Reconstructed audit text & Verbatim governing-specification text for two completion-report sections was unrecoverable after a mid-milestone context compaction; this paper reports the reconstructed audit, disclosed as such. \\
Single implementation & All results from one model lineage, one research team; independent re-implementation would strengthen confidence the continuum finding is not implementation-specific. \\
No historical-sequencing test & $H$ in $M_{\text{historical}}=M_{\text{general}}+H$ is not implemented or tested anywhere in this paper; it is a motivated hypothesis for future work only. \\
\bottomrule
\end{tabular}
}
\end{table}

\textbf{Parameter domain and duration.} Every result in this paper is
conditional on the tested mechanism set, the tested parameter domain, and
the tested run durations (up to 700 steps at Level 1.2b's production
scale). A \texttt{continuum} verdict is a statement about the outcome
space actually sampled; it is not a proof that no parameter region or
longer duration exists in which discrete institutional clustering would
emerge. \cref{sec:appendix-tables} reports the full sampled ranges.

\textbf{Collapse remains the modal outcome.} Even at Level 1.2b's much
larger and ecologically stratified scale, collapse remained the single
largest outcome class in every pooled sample (56.6\%--78.1\% depending on
the sample). Non-collapse institutional variation, while far more
plentiful than at Level 1.2 (161 viable and 89 fragmented runs out of 576
at discovery, versus 2 non-collapse runs out of 24 at Level 1.2), is still
a minority of the sampled space; a design deliberately oversampling the
non-collapse region --- informed by, but going beyond, the causal
ablation's own findings about subsistence cost and migration --- would
give any future discovery attempt more institutional variation to work
with, at the cost of no longer representing the model's own unconditional
outcome distribution.

\textbf{Evidence gaps disclosed, not filled.} \texttt{shared\_lineage\_%
persistence}, one of six candidate explanations Work Package O considered
for cross-run phenotype similarity, was never evaluated; doing so would
require a fourth full 576-run re-simulation with per-run lineage
identifiers exported, which the source project judged low-value given
that the continuum and null-replication findings already independently
established that no discrete candidate structure exists for shared
ancestry to explain. We report this as an open evidence gap rather than
filling it by rerunning simulations for this paper, per this paper's own
task discipline of not rerunning expensive computation merely to
strengthen a narrative.

\textbf{Temporal-representation validation subset.} Level 1.2b's four
burn-in-sensitivity definitions and window-bounds machinery were validated
on a disclosed 12-seed-run subset of the main design, not the full 576-run
pool, for compute-budget reasons stated at the time rather than after
seeing results.

\textbf{Supplementary statistical tests.} The causal-ablation Wilcoxon
$p$-values reported in \cref{sec:level12-collapse} (Level 1.2, $n=8$
matched pairs per test) are explicitly supplementary, not confirmatory ---
too small a sample for a standalone significance claim on their own, and
reported as such rather than overclaimed. Level 1.2b's 13-factor
decomposition (\cref{sec:level12b-scale-causal}) reports effect sizes and
outcome-label transition counts rather than significance tests, by design.

\textbf{Reconstructed audit sections.} Two portions of Level 1.2b's own
completion documentation --- the exact verbatim text of its governing
specification's stop-condition and final-report sections --- were not
recoverable from that project's own session context after a mid-milestone
context compaction, and its own completion report discloses this
explicitly rather than fabricating numbered items to match a remembered
count. This paper reports the reconstructed audit the source project
performed in its place, not a verbatim reproduction of text that no
longer exists.

\textbf{Single research team, single model lineage.} Every result in this
paper comes from one model implementation, developed and audited by one
project, however extensively internally cross-checked (structural
firewalls, byte-for-byte no-op regression tests, independent ledger
verification, disjoint confirmation samples). External replication in an
independently implemented ABM sharing the same generative assumptions
would strengthen confidence that the continuum finding reflects the
mechanism set itself rather than an implementation-specific artifact, in
the spirit of the broader empirical-validation standards recommended for
agent-based modeling generally \parencite{windrum2007empirical}.

\textbf{Historical-sequencing test now performed, with its own disclosed
scope.} Unlike the Level 1.2/1.2b manuscript this paper revises, $H$ is now
implemented and tested (\cref{sec:level13,sec:level13b}). Its own scope is
bounded: Stage 5's persistence diagnostics and Stage 6's history-erasure
interventions both operate at one representative cell
(\texttt{entrenchment=moderate}, \texttt{legitimacy=dual},
\texttt{stratum=medium\_feasibility}), not the full Stage 3 factorial;
Stage 6's own causal design manipulates historical state (pressure memory,
entrenchment age, legitimacy scores) rather than testing every conceivable
historical intervention; and every Level 1.3b stage ran at a disclosed,
session-budget-scaled \textit{minimum} sample size (Stage 4: 240 of a
preferred 432--768; Stage 8: 96 of a preferred 144--288), not adjusted
after seeing results.

\textbf{Discrete-versus-continuous comparison uses non-commensurable
metrics.} Level 1.3b's own Stage 8 compares the leading discrete candidate
(log-likelihood) against the leading continuous candidate (reconstruction
error) on fresh data --- two metrics not on a common scale
(\cref{sec:level13b-h7h8}), one of the decision gate's own four concrete
``what would change this verdict'' items (extending the tested $k$ range
past 8; building a common discrete/continuous fresh-data metric; extending
Stage 5/6 to the full Stage 3 factorial; building a genuine multi-condition
cluster-assignment-stability diagnostic for H3).

\textbf{$S_{\text{obs}}$/$Z_{\text{obs}}$ are not the dissertation's SEM
latent factors.} Every claim in \cref{sec:disc-relationship-to-dissertation}
treating the ABM's own observer-side composites as weak corroboration of
the dissertation's $S$/$Z$ framework is explicitly bounded to that weak
sense; $S_{\text{obs}}$/$Z_{\text{obs}}$ were never fit to, validated
against, or constructed with reference to the dissertation's own survey
data or SEM output, and the dissertation itself reports no single $S$--$Z$
correlation coefficient for this paper to compare against in the first
place (\cref{sec:appendix-ledger}).

\textbf{Reconstructed governing specifications.} Both Level 1.3 and Level
1.3b's own governing specifications were authored at length in an earlier
working session whose exact verbatim text became unrecoverable after a
later context compaction; every downstream document (preregistration
modules, scientific notes, decision-gate reports) discloses this
explicitly rather than presenting a reconstructed hypothesis list as a
verbatim transcription, and this paper reports the reconstructed
hypotheses and findings the source project itself worked from
(\cref{sec:reproducibility}).

\textbf{Real governance systems, not simulated ones, remain untested.}
Every claim in this paper about historical path dependence, causal
divergence, and continuous versus discrete structure is a claim about
\textit{this generative model's own outcome space}, under the specific
mechanism set, parameter domain, and duration tested. Whether real
governance systems' own historical trajectories are governed by dynamics
of the same qualitative shape --- and, in particular, why real systems
appear to concentrate around a small number of recurring configurations,
if the dissertation's own three-trajectory reading is correct --- is not
addressed by any simulation result in this paper and is deferred to
future, historically calibrated research (\cref{sec:disc-historical-sequencing}).

\Cref{tab:conclusion-boundary} states, for every major claim in this
paper, whether it is supported, partially supported, unsupported, or
reserved for future work, as a single strict reference table.

\begin{table}[htbp]
\centering
\caption{Strict conclusion-boundary table: every major claim in this
paper, classified by evidentiary status.}
\label{tab:conclusion-boundary}
\adjustbox{max width=\textwidth}{\begin{tabular}{@{}p{6.2cm}p{2.4cm}p{6.2cm}@{}}
\toprule
\textbf{Claim} & \textbf{Status} & \textbf{Basis / what would change it} \\
\midrule
Atomic mechanisms are generatively sufficient to produce institutions and heterogeneous institutional outcomes & Supported & Levels 1.0--1.2b; unforced formation at every tested ecological candidate \\
Synchronic substrate does not compress outcomes into discrete recurring forms & Supported & Level 1.2/1.2b: continuum, independently replicated (null\_replicated) \\
Historically structured pressure/sequencing causally shapes institutional trajectory & Supported & Level 1.3b Stage 6: sham-validated, 3/4 real interventions show growing divergence \\
Historically extended substrate compresses outcomes into discrete recurring forms & Unsupported & Level 1.3b: silhouette never clears 0.15; weak cross-algorithm/cross-stratum agreement \\
Pressure-onset ordering produces a real effect & Partially supported & Real but modest, direction-inconsistent across ordering pairs (H2) \\
Entrenchment moderates persistence as a graded dose-response & Unsupported (as stated) & Mechanism works; composite effect non-monotonic (H4) \\
$S_{\text{obs}}$--$Z_{\text{obs}}$ are coupled & Supported, modest & $r\approx0.26$--$0.29$ across variants; not the same construct; PCA-weighted variant $r\approx0.14$ \\
ABM results validate the dissertation's 3 empirical trajectories & Unsupported --- not attempted & Different analytical levels; ABM cannot confirm/refute a cross-sectional SEM typology \\
Real governance systems are discretely speciated into 3 (or any fixed count of) trajectories & Out of scope & Requires historically calibrated real-world data, not simulation \\
$\mathcal{G}_t$ richer-state formalism improves on $S$/$Z$ alone for tracking simulated trajectories & Supported, within this model & Motivated by Level 1.3/1.3b's own findings; not independently validated against external data \\
Why real systems concentrate around specific recurring trajectories & Deferred & Future, historically calibrated research (\cref{sec:disc-historical-sequencing}) \\
\bottomrule
\end{tabular}
}
\end{table}
